%\paragraph{Defining PVS}
To define the language of PVS in \mmt, we carry out two steps.

Firstly, we choose a suitable logical framework by picking the necessary features presented in \autoref{ch:lf}. Notably, we include three features: anonymous
record types (see \autoref{sec:lf:records}), predicate subtypes (see \autoref{sec:lf:lfx:subtp}), and imports of multiple instances of the same parametric
theory. The latter is achieved by a structural feature (see \autoref{sec:lf:lfx:wtypes}) that elaborates into the lambda-abstracted declarations of the imported theory.

Then we use this logical framework to define the \mmt theory for PVS. \autoref{fig:pvsmmtbasics} shows the most fundamental constants of this theory.

%\begin{figure}[ht]\centering
\begin{mmtcode}[caption={The Fundamental Higher-Order Logic of PVS \protect\footnotemark},label=fig:pvsmmtbasics]
tp : type &\RS&
expr : tp → type &\US& # 1 prec -1 &\RS&
tpjudg : {A} expr A → tp → type &\US& # 2 : 3 &\RS&

pvssigma : {A} (expr A ⟶ tp) ⟶ tp ❘ # Σ 2 ❙
tuple_tp : tp ⟶ tp ⟶ tp ❘ #  1 × 2  ❘ = [A][B] Σ [x:expr A]B ❙
tuple_expr : {A,B} expr A ⟶ expr B ⟶ expr A × B ❘ # < 3 , 4 > ❘ ## ( 3 , 4 ) ❙
proj : {A,B} expr A ⟶ expr NatLit ⟶ expr B ❘ # 3 _ 4 ❙

pvspi : {A} (expr A → tp) → tp &\US& # Π 2 &\RS&
fun_type : tp → tp → tp &\US& = [A,B] Π [x: expr A] B &\US& # 1 ⟹ 2 &\RS&
pvsapply : {A,f : expr A → tp} expr (Π f) → {a:expr A} expr (f a) 
	&\US& # 3 ( 4 ) prec -1000015 &\RS&
lambda : {A,f : expr A → tp} ({a:expr A}expr (f a)) → expr (Π f) &\US& # λ 3 &\RS&

bool : tp ❙
boolean : tp ❘ = bool ❙
  
FALSE : expr bool ❙
TRUE: expr bool ❙
NOT @ ¬ ❘ : expr (bool ⟹ bool) ❙
...

NatLit : tp ❙
StringLit : tp ❙
RatLit : tp ❙
rule rules?NatLiterals ❙
rule rules?StringLiterals ❙
rule rules?RationalLiterals ❙	
...

setsub : {A} (expr (A ⟹ bool)) ⟶ tp ❘ # ⟨ 1 | 2 ⟩ ❙
rule rules?SetsubRule ❙

rectp : {' '} ⟶ tp ❘ ## 1 ❙
recordexpr : {r : {' '} } {' '} ⟶ expr (rectp r) ❘ ## 2 ❙
fieldapp # 1 .@ 2 ❘  ## 1 . 2 ❙

recursor : {A} (expr A ⟶ expr A) ⟶ expr A ❘ ## INDUCTIVE 2❙

\end{mmtcode}
%\caption{Some Basic Typing related Symbols for PVS\protect\footnotemark}\label{fig:pvsmmtbasics}
%\end{figure}
\footnotetext{\url{https://gl.mathhub.info/PVS/Prelude/blob/master/source/pvs.mmt}}

We begin with a definition of PVS's higher-order logic using only LF features.  This
includes dependent product and function types\footnote{Contrary to typical
  dependently-typed languages, PVS does not allow declaring dependent \emph{base} types,
  but predicate subtyping can be used to introduce types that depend on
  terms. Interestingly, this is neither weaker nor stronger than the dependent types in
  typical $\lambda\Pi$ calculi.}, classical booleans, and the usual formula constructors
(see \autoref{fig:pvsmmtbasics}).  This is novel in how exactly it mirrors the syntax
of PVS (e.g., PVS allows multiple aliases for primitive constants) but requires no special
\mmt features.

We declare three constants for the three types of built-in literals together with \mmt
rules for parsing and typing them.  Using the new framework features, we give a shallow
encoding of predicate subtyping (see \autoref{fig:pvsmmtsubtypes} for the new typing
rule), a shallow definition of anonymous record types, as well as new declarations for
PVS-style inductive and co-inductive types.

%\begin{figure}[ht]\centering
%\hrule
\begin{scalacode}[caption={PVS-Style Predicate Subtyping in MMT and the Corresponding Rule \protect\footnotemark},label=fig:pvsmmtsubtypes]
object SetsubRule extends ComputationRule(PVSTheory.expr.path) {
  def apply(check: CheckingCallback)(tm: Term, covered: Boolean)
    (implicit stack: Stack, history: History): Option[Term]
   = tm match {
    case expr(PVSTheory.setsub(tp,prop)) => 
      Some(LFX.Subtyping.predsubtp(expr(tp),proof("internal_judgment",
        Lambda(doName,expr(tp),pvsapply(prop,OMV(doName),expr(tp),
        		bool.term)._1))))
    case _ => None
  }
}
\end{scalacode}
%\caption{PVS-Style Predicate Subtyping in MMT and the Corresponding Rule \protect\footnotemark}\label{fig:pvsmmtsubtypes}
%\end{figure}
\footnotetext{Ibid. and \url{https://github.com/UniFormal/MMT/blob/master/src/mmt-pvs/src/info/kwarc/mmt/pvs/Plugin.scala}}